% ==================================================================
% id:       rv-001
% topic:    משתנים מקריים
% subtopic: פונקציית הסתברות ופונקציית התפלגות מצטברת
% solution: full
% verified: no
% ==================================================================
נתון משתנה מקרי $X$ עם טווח ערכים $R_X=\set{0,1,2,3,4}$, וידוע ש:
\begin{itemize}
\item $F_X(2)=F_X(3)=0.8$;
\item $\PP(X=0)+\PP(X\ge2)=0.5$;
\item $\PP(X>0)-\PP(X<3)=0.1$.
\end{itemize}
מצאו את פונקציית ההסתברות $P_X(k)$ לכל $k\in R_X$.
\hint{סמנו $p_k=\PP(X=k)$ וכתבו כל אחד מהנתונים כמשוואה, למשל $F_X(2)=p_0+p_1+p_2$. אל תשכחו ש-$p_0+p_1+p_2+p_3+p_4=1$.}
\begin{solbox}
\unverified
נפענח את הנתונים אחד-אחד, בעזרת העובדה ש-$R_X=\set{0,1,2,3,4}$.
\begin{itemize}
\item $F_X(3)=\PP(X\le2)+\PP(X=3)=F_X(2)+\PP(X=3)$, ולכן מהנתון הראשון $\ans{\PP(X=3)=0}$.
\item גם מהנתון הראשון: $\PP(X=4)=\PP(X>3)=1-F_X(3)=\ans{0.2}$.
\item מהנתון השני: $0.5=\PP(X=0)+\PP(X\ge2)=1-\PP(X=1)$, ולכן $\ans{\PP(X=1)=0.5}$.
\item מהנתון השלישי: $\PP(X>0)=1-\PP(X=0)$ ו-$\PP(X<3)=\PP(X\le2)=0.8$, ולכן $1-\PP(X=0)-0.8=0.1$, כלומר $\ans{\PP(X=0)=0.1}$.
\item לסיום, $\PP(X=2)=1-(0.1+0.5+0+0.2)=\ans{0.2}$.
\end{itemize}
\end{solbox}
